\section*{Scene, composition and affordance}
Warm off-white walls, charcoal worktops, muted blue-green storage, metal supports and restrained timber edging suggest a contemporary teaching laboratory. Ceiling diffusers, ventilation, an element poster, reagent archive, safety signage, notebooks and a distant instrument workstation provide context without competing with the assembly task.

\begin{center}
\begin{tikzpicture}[x=1.25cm,y=1.05cm,font=\small]
\fill[pale] (-3,-2.5) rectangle (3,2.5);\draw[thick] (-3,-2.5) rectangle (3,2.5);
\fill[accent!45] (-2.85,1.8) rectangle (2.85,2.4);
\node at (0,2.1) {Fixed exhibits / instruments};
\draw[fill=ink!25] (-1.025,-.425) rectangle (1.025,.425);
\node[font=\scriptsize] at (-.24,.14) {Assembly};
\draw[fill=accent!25] (.45,-.3) rectangle (.98,.05);\node[font=\tiny] at (.71,-.13){Tray};
\draw[fill=accent!35] (-1.38,-1.45) rectangle (-.84,-.48);
\draw[<-] (-1.38,-.7)--(-1.8,-.7) node[left,font=\scriptsize]{Shelf};
\draw[fill=white] (0,-.85) circle (.16);\node[below,font=\scriptsize] at (0,-1.08){Learner};
\draw[dashed,accent,thick] (0,-.85) circle (.85);
\draw[<-] (.73,-1.3)--(1.45,-1.65) node[right,align=left,font=\scriptsize]{Intended\\task reach};
\draw[<->] (-3,2.8)--node[above]{6 m}(3,2.8);
\draw[<->] (3.35,-2.5)--node[right]{5 m}(3.35,2.5);
\end{tikzpicture}
\par\small Figure 1. Design floor plan; dimensions describe the virtual room.
\end{center}


The table remains central in the room. The learner occupies its near edge; stock is beside the left hand and the reaction tray is on the right. Shelf and bench height move together; forward/back adjustment supports reach. The shelf has moved 20 cm left and forward to provide 5 cm worktop clearance. The diagram shows intended reach relationships, not a certified ergonomic envelope. Seated and standing reach must be tested on Quest.

\begin{tabularx}{\linewidth}{@{}p{31mm}YY@{}}\toprule
Object class & Visual language & Behaviour\\\midrule
Atoms and tools & Element symbols, clear silhouettes, hover feedback & Respond to grip/pinch; provide immediate state feedback\\
Fixed shelf/baskets & Matte tray surfaces, sockets and count labels & Stay fixed; only contents can be picked up\\
Background equipment & Muted materials; positioned beyond task reach & No hover response or grab interaction\\
Bond targets & Direction guides, arc and release ring & Appear only for valid authored opportunities\\\bottomrule
\end{tabularx}

Avoid nearby nonfunctional handles or convincing grab cues. Physical placement, material treatment and behavioural consistency must reinforce each other; colour alone is insufficient. Background objects use inexpensive opaque materials, including simplified glassware.
